% ===========================================================================
%  Resumo, abstract
% ===========================================================================
\chapter*{RESUMO}
\addcontentsline{toc}{chapter}{RESUMO}
\thispagestyle{plain}

\SingleSpacing
\noindent
Este trabalho descreve a constru\c{c}\~ao de um modelo tridimensional do autom\'ovel esportivo
Lotus Elise 111S (Series 1) integralmente por c\'odigo, sem utiliza\c{c}\~ao de malhas de
terceiros, e a integra\c{c}\~ao desse processo a um agente de intelig\^encia artificial por meio do
\emph{Model Context Protocol} (MCP). O problema abordado \'e a aus\^encia de um pipeline
reprodut\'ivel que una especifica\c{c}\~ao textual, refer\^encias fotogr\'aficas e edi\c{c}\~ao
tridimensional interativa: na pr\'atica industrial, ajustes de malha tendem a se perder em
sess\~oes manuais n\~ao versionadas. A pesquisa adota abordagem de engenharia reversa visual:
um conjunto de seis fotografias do ve\'iculo real \'e tratado como \emph{ground truth}; dessas
refer\^encias derivam-se dimens\~oes de envelope, esta\c{c}\~oes de casco e par\^ametros de
materiais, formalizados em especifica\c{c}\~ao pr\'opria segundo o m\'etodo de Desenvolvimento
Orientado por Especifica\c{c}\~ao. A geometria \'e gerada por um \emph{loft} de 45 esta\c{c}\~oes
transversais com fundo modulado para produzir arcos de roda reais sem opera\c{c}\~oes booleanas,
garantindo malha quadr\^angular dominante e aus\^encia de n-gons. Os materiais empregam o
\emph{Principled BSDF} com pintura automotiva multicamada, \emph{clear coat} e micro-flocos
met\'alicos. A verifica\c{c}\~ao \'e visual e comparativa, com registro de aceite por crit\'erio.
Os resultados obtidos totalizam 130 objetos de malha, 32.632 faces, zero n-gons e zero arestas
n\~ao variedade, com envelope de 3,736\,m $\times$ 1,718\,m $\times$ 1,117\,m, contra alvos de
3,726\,m, 1,718\,m e 1,117\,m. Nove renders de verifica\c{c}\~ao e cinco imagens de galeria foram
produzidos no motor EEVEE. Dos nove crit\'erios de aceita\c{c}\~ao, sete foram considerados
plenamente atendidos e dois, parcialmente: a silhueta carece da linha de car\'ater lateral viva e
o tom da pintura ficou ligeiramente mais escuro que a refer\^encia. Conclui-se que a combina\c{c}\~ao
entre especifica\c{c}\~ao formal, geometria procedimental e automa\c{c}\~ao via MCP produz um
artefato audit\'avel, reconstru\'ivel por um \'unico comando e adequado \`a verifica\c{c}\~ao
independente.

\vspace{1.0cm}
\noindent\textbf{Palavras-chave}: Model Context Protocol; Blender; modelagem procedimental;
subdivis\~ao de superf\'icies; renderiza\c{c}\~ao fisicamente baseada; engenharia reversa visual.

\clearpage

% ===========================================================================
\chapter*{ABSTRACT}
\addcontentsline{toc}{chapter}{ABSTRACT}
\thispagestyle{plain}

\SingleSpacing
\noindent
This work describes the construction of a three-dimensional model of the Lotus Elise 111S
(Series 1) sports car entirely through code, without third-party meshes, and the integration of
that process with an artificial-intelligence agent through the \emph{Model Context Protocol}
(MCP). The problem addressed is the lack of a reproducible pipeline that unifies textual
specification, photographic references and interactive three-dimensional editing: in industrial
practice, mesh adjustments tend to be lost in unversioned manual sessions. The research adopts a
visual reverse-engineering approach: a set of six photographs of the real vehicle is treated as
ground truth; envelope dimensions, hull stations and material parameters are derived from these
references and formalised in a specification following the Specification-Driven Development
method. Geometry is generated by a 45-station cross-section loft whose bottom is modulated to
produce real wheel arches without boolean operations, ensuring a quad-dominant mesh with no
n-gons. Materials employ the \emph{Principled BSDF} with multi-layer automotive paint, clear coat
and metallic micro-flakes. Verification is visual and comparative, with acceptance recorded per
criterion. The results total 130 mesh objects, 32,632 faces, zero n-gons and zero non-manifold
edges, with an envelope of 3.736\,m $\times$ 1.718\,m $\times$ 1.117\,m against targets of
3.726\,m, 1.718\,m and 1.117\,m. Nine verification renders and five gallery images were produced
with the EEVEE engine. Of the nine acceptance criteria, seven were fully met and two partially:
the silhouette lacks the sharp character line and the paint tone came out slightly darker than
the reference. It is concluded that the combination of formal specification, procedural geometry
and MCP-driven automation yields an auditable artefact, reconstructible with a single command and
suitable for independent verification.

\vspace{1.0cm}
\noindent\textbf{Keywords}: Model Context Protocol; Blender; procedural modelling; surface
subdivision; physically based rendering; visual reverse engineering.
\OnehalfSpacing
\clearpage
